\section{Introduction}\label{sec:introduction}

\subsection{Context and significance}

Deep neural networks now provide highly effective representations for image
classification, detection, segmentation, and many downstream vision tasks.
Their predictive success, however, does not by itself reveal how evidence is
formed inside the model. A classification score is the endpoint of a sequence
of spatial transformations: local edges and textures are combined into parts,
parts are grouped into objects, and object representations become sufficiently
class-discriminative for the final decision. Understanding this progression is
important for at least three reasons. First, model developers need diagnostic
tools that reveal failure modes rather than only display a final confidence
score. Second, users need explanations whose meaning and limitations are
explicit. Third, internal structure that is genuinely useful should support
interventions such as model compression or knowledge transfer.

Visual explanation research has largely approached this problem through
attribution. Class activation mapping (CAM) and its descendants project
class-relevant signals onto a spatial feature map
\cite{zhou2016cam,selvaraju2017gradcam,chattopadhay2018gradcampp}. These
methods have made model decisions easier to inspect, but a heatmap is not a
complete account of representation formation. It normally begins with a chosen
output and a chosen layer, then asks which spatial locations support that
output. It does not directly identify where the internal representation changes
from generic visual encoding to mature class organization. This project studies
that earlier and broader question.

\subsection{Motivation}

The work began with dataset-wide principal component analysis (PCA) of
intermediate convolutional features. For each layer, a basis is fitted using
one dataset-wide pool of spatial vectors sampled across the training split. The
first three coordinates are rendered as red, green, and blue channels, allowing
within-layer comparisons across images because all images share the same
layer-specific basis. The resulting maps suggested a recurring qualitative
progression: early maps preserve local image structure; intermediate maps
undergo conspicuous spatial reorganization; and late maps become more
class-oriented.

Several observations made this progression scientifically interesting. A
randomly initialized CNN can retain visible image structure at unexpectedly
late depth, indicating that architecture alone can propagate structured
signals. Training changes that structure into a more class-discriminative
organization. On CIFAR-10, images from the same class can show similar
late-layer PCA-RGB organization, and selected misclassified images appear more
similar to the predicted class than to the labelled class. On CIFAR-100 and
ImageNet-100, such class-consistent colours are less visually obvious at late
depth, but intermediate reorganization remains observable. These findings
motivate a transition-layer hypothesis, while also demonstrating why visual
inspection alone is insufficient.

The hypothesis led to two methodological directions. The first, \methodname,
asks whether the geometry of late unpooled spatial features can itself provide
an explanation. It models recurrent dataset-global feature modes using robust
prototypes and scores each spatial position by its distance from the nearest
mode. The second, \jmethod{}---the project name for a selective hierarchical
knowledge-distillation objective---asks whether selected spatial and relational
geometry can be transferred from a teacher to a smaller student. This creates a
linked programme in which representation geometry is first observed, then used
for explanation, and finally tested as an actionable training signal.

\subsection{Problem statement}

The project addresses four related problems.

\begin{enumerate}[leftmargin=*]
    \item Existing CAM methods provide useful output-conditioned localization,
    but do not fully characterize how spatial representations evolve across
    depth or how an explanatory layer should be selected.
    \item PCA-RGB can reveal repeatable visual patterns, but colour similarity
    is not automatically semantic evidence. Each layer has its own basis, so
    direct colour identity across layers is not defined, and global variance may
    reflect texture, background, position, or scale.
    \item Competitive object localization does not automatically imply
    classifier faithfulness. A class-agnostic representation map and a
    class-conditioned attribution map answer different questions and require
    different evaluation.
    \item Knowledge-distillation methods often match logits, features, or
    relations without an explicit account of which stage of representation
    formation should be constrained. Excessive layer-wise imitation can
    overconstrain a compact student whose computational path differs from its
    teacher.
\end{enumerate}

The central scientific problem is therefore to identify a reproducible
representation-transition regime, determine what explanatory information its
geometry contains, and test whether that information is transferable.

\subsection{Aim, research questions, and objectives}

The overall aim is:

\begin{quote}
To understand how spatial representations become class-discriminative across
the depth of vision models, and to use this understanding to develop visual
explanation and knowledge-transfer methods.
\end{quote}

Four research questions organize the programme:

\begin{description}[leftmargin=1.3cm,style=nextline]
    \item[RQ1] How do learned spatial representations evolve across CNN depth,
    and can the transition toward class-discriminative organization be detected
    reproducibly?
    \item[RQ2] Can transition-layer geometry provide informative and faithful
    visual explanations?
    \item[RQ3] Can related geometry be transferred from teacher to student to
    improve learning, and which layers and geometric components are responsible?
    \item[RQ4] Do comparable transitions arise in token-based vision models,
    and what explanation or transfer opportunities do they enable?
\end{description}

The corresponding objectives are to: (i) construct aligned representation
visualizations and quantitative transition measures; (ii) develop and evaluate
\methodname{} as a class-agnostic explanation method; (iii) develop and audit
the \jmethod{} geometry-distillation objective; and (iv) investigate
transformer representations without prematurely fixing one architecture or
method.

\subsection{Expected contributions and scope}

The expected contributions are:

\begin{enumerate}[leftmargin=*]
    \item an empirical and quantitative characterization of spatial
    representation transition across trained CNNs;
    \item a class-agnostic, gradient-free explanation principle based on
    distance from adaptive dataset-global robust feature modes;
    \item a selective hierarchical distillation method that transfers early
    relations, balanced transition coordinates, and late relational evolution;
    \item an evaluation framework that distinguishes object localization,
    class attribution, representation diagnostics, and transferable utility;
    and
    \item evidence clarifying whether the same transition hypothesis extends
    beyond CNN spatial grids to transformer token representations.
\end{enumerate}

The current evidence primarily concerns image classification and CNN families.
The title uses the broader term vision models as the intended thesis scope, but
claims in the completed-work chapters are restricted to the architectures and
datasets actually tested. Transformer generalization remains prospective.

\subsection{Research programme}

\begin{figure}[H]
    \centering
    \begin{tikzpicture}[
    node distance=7mm,
    stage/.style={draw,rounded corners,minimum width=0.78\linewidth,
                  minimum height=10mm,align=center,font=\small},
    arrow/.style={-{Latex[length=2.3mm]},thick}
]
\node[stage,fill=stageblue] (observe)
    {Observe: dataset-wide PCA-RGB maps across model depth};
\node[stage,fill=stageblue,below=of observe] (validate)
    {Validate: controls, quantitative transition measures, and failure analysis};
\node[stage,fill=stagegreen,below=of validate] (explain)
    {Explain: class-agnostic robust spatial geometry with Adaptive K-DGRP};
\node[stage,fill=stageorange,below=of explain] (transfer)
    {Transfer: selective hierarchical geometry distillation with J1};
\node[stage,fill=gray!12,below=of transfer] (generalise)
    {Generalize: test the hypothesis in token-based vision models};
\draw[arrow] (observe) -- (validate);
\draw[arrow] (validate) -- (explain);
\draw[arrow] (explain) -- (transfer);
\draw[arrow] (transfer) -- (generalise);
\end{tikzpicture}

    \caption{Research programme connecting observation, controlled validation,
    explanation, knowledge transfer, and future generalization.}
    \label{fig:research-program}
\end{figure}

\Cref{fig:research-program} shows the intended logic. PCA-RGB is a discovery
and diagnostic instrument, not by itself a faithful attribution method.
\methodname{} tests whether recurrent feature-space geometry can localize
object-relevant structure. \jmethod{} tests whether selected geometry has
training utility. These studies are complementary, but the report avoids the
stronger unsupported claim that successful distillation proves the faithfulness
of a particular explanation map.
